% =====================================================================
%  1.8 Implementação dos sistemas de esgotamento sanitário
%  (Atividades 3.2.2.4 e 3.2.2.5)
% =====================================================================
\section{Implementação dos sistemas de esgotamento sanitário}
\label{sec:esgotamento}

O esgotamento será feito com soluções individuais em cada domicílio, em duas
frentes: 4.067 domicílios nas comunidades atendidas com água (2028 a 2033) e
2.300 domicílios em mananciais, na área do IDR-Paraná (2029 a 2033)
\cite{custo2026}. As frentes diferem na origem da demanda, mas seguem o mesmo
procedimento a partir do projeto de cada propriedade
(\autoref{fig:esgotamento-geral}).

\begin{figure}[htbp]
  \centering
  \caption{Implementação do esgotamento sanitário, do começo ao fim}
  \label{fig:esgotamento-geral}
  \fluxograma{fluxograma-esgotamento-geral}
  \fonte{Elaborado pelo IAT (2026) com base em \citeonline{mop2026} e no POA 2026--2027.}
\end{figure}

\pendencia{O MOP não indica a tecnologia de tratamento nem as normas de
projeto. Após o estudo de alternativas (POA, out--dez/2026), citar as normas
aplicáveis, como a ABNT NBR 7229 e a ABNT NBR 13969. O POA fala em ``rede de
esgotamento'', termo que não corresponde a soluções individuais.}

% ---------------------------------------------------------------------
\subsection{Instalação e verificação técnica}
\label{subsec:esgotamento-instalacao}

A empresa contratada projeta e instala o sistema em cada propriedade e
capacita o proprietário durante a instalação. O sistema só segue para medição
depois de aprovado na verificação técnica. Se reprovado, a empresa corrige e o
sistema é verificado de novo (\autoref{fig:esgotamento-instalacao}).

\begin{figure}[htbp]
  \centering
  \caption{Projeto, instalação e verificação técnica dos sistemas individuais}
  \label{fig:esgotamento-instalacao}
  \fluxograma{fluxograma-esgotamento}
  \fonte{Elaborado pelo IAT (2026) com base em \citeonline{mop2026}.}
\end{figure}

% ---------------------------------------------------------------------
\subsection{Medição, comprovação e pagamento}
\label{subsec:esgotamento-medicao}

O pagamento depende de uma cadeia de comprovação. A empresa mede e documenta,
o proprietário assina a confirmação, e o IAT faz auditoria documental,
verificação e amostragem \textit{in loco} antes de autorizar o pagamento
(\autoref{fig:esgotamento-medicao}).

\begin{figure}[htbp]
  \centering
  \caption{Medição, comprovação, verificação e pagamento, por responsável}
  \label{fig:esgotamento-medicao}
  \fluxograma{fluxograma-esgotamento-medicao}
  \fonte{Elaborado pelo IAT (2026) com base em \citeonline{mop2026}.}
\end{figure}

\pendencia{Definir o tamanho da amostra da verificação \textit{in loco} e
quem executa cada verificação.}

% ---------------------------------------------------------------------
\subsection{Frente dos mananciais (IDR-Paraná)}
\label{subsec:esgotamento-idr}

O \autoref{qua:diferencas-esgotamento} mostra o que muda entre as duas frentes.

\begin{quadro}[htbp]
\caption{Diferenças entre as frentes de esgotamento sanitário}
\label{qua:diferencas-esgotamento}
\small
\renewcommand{\arraystretch}{1.25}
\begin{tabularx}{\textwidth}{|P{3.2cm}|L|L|}
\hline
\cab{Aspecto} & \cab{Frente comunidades (3.2.2.4)} & \cab{Frente mananciais -- IDR-Paraná (3.2.2.5)} \\ \hline
Origem da demanda & Diagnóstico do trabalho técnico-social &
  Planos Integrados de Propriedade (PIP) do IDR-Paraná \\ \hline
Território & Comunidades atendidas com abastecimento de água &
  Mananciais prioritários do PSH-PR \\ \hline
Meta & 4.067 domicílios & 2.300 domicílios \\ \hline
Prazo & 2028 a 2033 & 2029 a 2033 \\ \hline
Dependências & Prefeituras & -- \\ \hline
Validação & IAT & IAT e IDR-Paraná \\ \hline
\end{tabularx}
\fonte{Adaptado de \citeonline{mop2026} e \citeonline{custo2026}.}
\end{quadro}
